\subsection{连续核积分算子的迹}

在本节中，我们保持前一节的约定，取 Y = X 且 \(\nu = \mu\)。我们假设 X 是在无穷远处可数的局部紧拓扑空间（TG, I, p. 68，定义 5）。

特别地，我们把空间 \(\mathrm{L}^{2}(\mathrm{X}\times\mathrm{X})\) 与 \(\mathrm{L}^{2}(\mathrm{X})\widehat{\otimes}_{2}\mathrm{L}^{2}(\mathrm{X})\) 等同起来（IV, p. 172 的引理 10）。我们用 \(\widetilde{f}\) 表示 \(\mathrm{L}^{2}(\mathrm{X})\) （相应地，\(\mathrm{L}^{2}(\mathrm{X}\times\mathrm{X})\) ）中函数 \(f\in\mathcal{L}^{2}(\mathrm{X})\) （相应地，\(\mathcal{L}^{2}(\mathrm{X}\times\mathrm{X})\) 中的函数）的类。

命题 20. --- 设 \((f_n)_{n \in \mathbf{N}}\) 是从 X 到可度量化空间 F 中的可测映射序列。假设 \(f_n(x)\) 的极限在 X 的一个 \(\mu\)-零测子集之外存在。则存在 X 的紧子集序列 \((\mathrm{C}_m)_{m \in \mathbf{N}}\) ，其并 \(\widetilde{\mathrm{X}}\) 满足 \(|\mu|(\mathrm{X} - \widetilde{\mathrm{X}}) = 0\) ，使得函数 \(f_n\) 在 \(\mathrm{C}_m\) 上连续（对所有 \(n \in \mathbf{N}\) 与每个 \(m \in \mathbf{N}\)），且序列 \((f_n)_{n \in \mathbf{N}}\) 一致收敛到 \(f\) （在 \(\mathrm{C}_m\) 上，对所有 \(m \in \mathbf{N}\)）。

由 INT, IV, p. 175, § 5, no. 4 的定理 2 及 INT, IV, p. 188, § 5, no. 8 的命题 12, b)，X 中使得函数 \(f_n\) 对每个 \(n\) 在 C 上连续且 \((f_n)\) 在 C 上一致收敛到 \(f\) 的紧子集 C 所成之集在 X 中 \(\mu\)-稠密（INT, IV, p. 189, § 5, no. 8，定义 6）。断言随即由 INT, IV, p. 189, § 5, no. 8 的注得出。

设 \(N\) 是属于 \(\mathcal{L}^2 (\mathrm{X}\times \mathrm{X})\) 的函数，\(u_{\mathrm{N}}\) 是从 \(\mathrm{L}^2 (\mathrm{X})\) 到 \(\mathrm{L}^2 (\mathrm{X})\) 中的以 \(N\) 为核的希尔伯特--施密特映射（IV, p. 173 的命题 19）。由于映射 \(u_{\mathrm{N}}\) 是紧的，由 IV, p. 156 的命题 9，存在元素 \(M\) （\(\overline{\mathbf{N}}\) 中的元素），并且用 \(I\) 表示整数 \(n\in \mathbf{N}\) 中满足 \(n\leqslant \mathrm{M}\) 者所成之集，存在标准正交族 \((f_i)_{i\in \mathrm{I}}\) 与 \((g_{i})_{i\in \mathrm{I}}\) （\(\mathcal{L}^2 (\mathrm{X})\) 中的标准正交族）以及族 \((\alpha_{i})_{i\in \mathrm{I}}\) （\(\mathbf{R}_+^*\) 中的族），使得

\[
u _ {\mathrm{N}} (f) = \sum_ {i \in \mathrm{I}} \alpha_ {i} \langle \widetilde {f} _ {i} | f \rangle \widetilde {g} _ {i}\tag{13}
\]

对所有 \(f \in \mathrm{L}^2(\mathrm{X})\) 成立，其中级数在 \(\mathrm{L}^2(\mathrm{X})\) 中收敛。

对每个 \(i \in I\) ，用 \(h_i \in \mathcal{L}^2(\mathrm{X} \times \mathrm{X})\) 表示由 \(h_i(x, y) = \overline{f_i(x)} g_i(y)\) 定义的函数，从而 \(\widetilde{h}_i\) 是 \(\overline{f}_i \otimes g_i\) 的类。由 IV, p. 173 的注 9，通项为 \(\alpha_i h_i\) 的级数在 \(\mathrm{L}^2(\mathrm{X} \times \mathrm{X})\) 中收敛到 N。

在本节其余部分，我们假设 N 是连续函数。

命题 21. --- 假设 \(u_{\mathrm{N}}\) 有有限迹。则存在一个集 \(\widetilde{\mathrm{X}} \subset \mathrm{X}\) ，其余集 \(\mathrm{X} - \widetilde{\mathrm{X}}\) 是 \(\mu\)-零测的，还存在函数 \(\mathrm{H} \in \mathcal{L}^2(\mathrm{X} \times \mathrm{X})\) ，满足以下条件：

\begin{enumerate}
\def\labelenumi{(\roman{enumi})}
\item
  对每个 \((x,y)\in\widetilde{\mathrm{X}}\times\widetilde{\mathrm{X}}\) ，族 \((\alpha_{i}\overline{f_{i}(x)}g_{i}(y))_{i\in\mathrm{I}}\) 在 \(\mathbf{C}\) 中可和，其和为 \(\mathrm{N}(x,y)\) ；
\item
  对 I 的每个有限子集 J 与每个 \((x,y)\in\widetilde{\mathbf{X}}\times\widetilde{\mathbf{X}}\) ，我们有
\end{enumerate}

\[
\left| \sum_ {i \in \mathrm{J}} \alpha_ {i} h _ {i} (x, y) \right| \leqslant \mathrm{H} (x, y);\tag{14}
\]

\begin{enumerate}
\def\labelenumi{(\roman{enumi})}
\setcounter{enumi}{2}
\tightlist
\item
  函数 \(x \mapsto \mathrm{H}(x, x)\) 属于 \(\mathcal{L}^1(\mathrm{X})\) 。
\end{enumerate}

由于 \(u_{N}\) 有有限迹，族 \((\alpha_{i})\) 可和（IV, p. 165 的推论 4）。级数

\[
\sum_ {i \in \mathrm{I}} \alpha_ {i} | f _ {i} | ^ {2}, \quad \sum_ {i \in \mathrm{I}} \alpha_ {i} | g _ {i} | ^ {2}
\]

因此在 \(\mathcal{L}^{1}(X)\) 中收敛，从而 \(\mu\)-几乎处处收敛。用 F 与 G 分别表示由这些级数的和在 \(\mu\)-几乎处处意义下定义的函数，其中置 \(F(x) = 0\) 与 \(G(x) = 0\) ，对相应级数不收敛的每个 \(x\) 。由于 F 与 G 属于 \(\mathcal{L}^{1}(X)\)，由 \(H(x,y) = \sqrt{F(x)G(y)}\) 定义的函数 H 属于 \(\mathcal{L}^{2}(X \times X)\) （INT, V, p. 95, § 8, no. 3，推论 2）。

把命题 20 应用于空间 X，取 \(F = \mathbf{C}^{2}\) ，并取可测映射

\[
s_{n}\colon x\mapsto \Bigl (\sum_{\substack{i\in \mathrm{I}\\ i\leqslant n}}\alpha_{i}|f_{i}(x)|^{2},\sum_{\substack{i\in \mathrm{I}\\ i\leqslant n}}\alpha_{i}|g_{i}(x)|^{2}\Bigr)
\]

它们定义在 X 上。因此存在 X 的紧子集序列 \((\mathrm{C}_{m})_{m\in\mathbf{N}}\) ，满足以下条件：

\begin{enumerate}
\def\labelenumi{(\arabic{enumi})}
\item
  并 \(\widetilde{X}\) 满足 \(\mu(\mathrm{X}-\widetilde{\mathrm{X}})=0\) ；
\item
  对所有 \(m \in \mathbf{N}\) 与每个 \(n \in \mathbf{N}\) ，函数 \(s_n\) 在 \(\mathrm{C}_m\) 上连续；
\item
  对所有 \(m \in \mathbf{N}\) ，级数 \(\sum \alpha_{i}|f_{i}|^{2}\) 与 \(\sum \alpha_{i}|g_{i}|^{2}\) 在 \(C_{m}\) 上分别一致收敛到 F 与 G；
\item
  由 \(\mu\) 在 \(C_{m}\) 上诱导的测度的支撑等于 \(C_{m}\) （必要时把 \(C_{m}\) 换成这个测度的支撑）。
\end{enumerate}

我们来证明集 \(\widetilde{\mathbf{X}}\) 与函数 H 满足条件 (i)、(ii) 与 (iii)。

设 \((x,y)\in \widetilde{\mathbf{X}}\times \widetilde{\mathbf{X}}\)。对每个有限子集 \(\mathrm{J}\subset \mathrm{I}\) ，我们有

\[
\left| \sum_ {i \in \mathrm{J}} \alpha_ {i} \overline {{f _ {i} (x)}} g _ {i} (y) \right| \leqslant \left(\sum_ {i \in \mathrm{J}} \alpha_ {i} | f _ {i} (x) | ^ {2}\right) ^ {1 / 2} \left(\sum_ {i \in \mathrm{J}} \alpha_ {i} | g _ {i} (y) | ^ {2}\right) ^ {1 / 2}\tag{15}
\]

进而有

\[
\left| \sum_ {i \in \mathrm{J}} \alpha_ {i} \overline {{f _ {i} (x)}} g _ {i} (y) \right| \leqslant \mathrm{H} (x, y),\tag{16}
\]

这已给出性质 (ii)。

由不等式 (15)，级数 \(\sum_{i} \alpha_{i} h_{i}\) 在 \(\mathrm{K}_{m} \times \mathrm{K}_{n}\) 上一致收敛（对所有 \((m, n) \in \mathbf{N}^{2}\)）。设 \(\widetilde{\mathrm{N}}\) 是 \(\mathrm{X} \times \mathrm{X}\) 上由

\[
\widetilde {\mathrm{N}} (x, y) = \sum_ {i \in \mathrm{I}} \alpha_ {i} h _ {i} (x, y) = \sum_ {i \in \mathrm{I}} \alpha_ {i} \overline {{f _ {i} (x)}} g _ {i} (y)
\]

对所有 \((x,y)\in\widetilde{\mathrm{X}}\times\widetilde{\mathrm{X}}\) 定义、否则置 \(\widetilde{\mathrm{N}}(x,y)=0\) 的函数。函数 \(\widetilde{N}\) 可测（INT, IV, p. 175, § 5, n°4，定理 2），且在 \(K_{m}\times K_{n}\) 上连续（对所有 \((m,n)\in\mathbf{N}^{2}\)）。

由 (16)，我们有 \(|\widetilde{\mathrm{N}}(x,y)| \leqslant \mathrm{H}(x,y)\) ，对所有 \((x,y) \in \mathrm{X} \times \mathrm{X}\) 成立，特别地 \(\widetilde{\mathrm{N}}\) 属于 \(\mathcal{L}^2(\mathrm{X} \times \mathrm{X})\) （INT, IV, p. 84, § 5, no. 6，定理 5）。

我们来证明 \(N = \widetilde{N}\) 在 \(\widetilde{X} \times \widetilde{X}\) 上成立，这将给出性质 (i)。设 f 与 g 是 \(\mathcal{L}^{2}(X)\) 的元素。我们有

\[
\langle g \mid u _ {\widetilde {N}} (f) \rangle = \int_ {X \times X} \widetilde {N} (x, y) f (x) \overline {{g (y)}} d (\mu \otimes \mu) (x, y)
\]

（IV, p. 172 的公式 (12)）。对每个 \((x,y)\in \widetilde{\mathrm{X}}\times \widetilde{\mathrm{X}}\) 与 I 的每个有限子集 J，我们有

\[
\left| \sum_ {i \in J} \alpha_ {i} \overline {{f _ {i} (x)}} g _ {i} (y) f (x) \overline {{g (y)}} \right| \leqslant | f (x) g (y) | H (x, y)
\]

（公式 (16)）。由于这个不等式的右端在 X × X 上可积（INT, V, p. 95, § 8, n° 3，推论 2），可以应用勒贝格定理（INT, IV, p. 137, § 3, n° 7，定理 6）及公式 (13)，推出

\[
\begin{array}{c} \langle g | u _ {\widetilde {N}} (f) \rangle = \sum_ {i \in I} \alpha_ {i} \int_ {X \times X} \overline {{f _ {i} (x)}} g _ {i} (y) f (x) \overline {{g (y)}} d (\mu \otimes \mu) (x, y) \\ = \sum_ {i \in I} \alpha_ {i} \langle f _ {i} | f \rangle \langle g | g _ {i} \rangle = \langle g | u _ {N} (f) \rangle . \end{array}
\]

因此可得 \(u_{\widetilde{\mathrm{N}}} = u_{\mathrm{N}}\) ，从而 \(\mathrm{N} = \widetilde{\mathrm{N}}\) 在 \(\mathrm{L}^2(\mathrm{X} \times \mathrm{X})\) 中成立 （III, p. 28 的命题 3, b)））。

对每个 \((m,n)\in\mathbf{N}^{2}\)，函数 N 与 \(\widetilde{N}\) 在 \(K_{m}\times K_{n}\) 上连续，故在 \(K_{m}\times K_{n}\) 上相等（INT, III, p. 69, § 2, n°2 的命题 9），因为 \(\mu\otimes\mu\) 在 \(K_{m}\times K_{n}\) 上诱导的测度的支撑等于 \(K_{m}\times K_{n}\) 。因此 N 与 \(\widetilde{N}\) 在 \(\widetilde{X}\times\widetilde{X}\) 上相合。

最后，上界 \(\sqrt{\mathrm{FG}} \leqslant (\mathrm{F} + \mathrm{G}) / 2\) 蕴含函数 \(x \mapsto \sqrt{\mathrm{F}(x)\mathrm{G}(x)} = \mathrm{H}(x,x)\) 在 X 上可积，由此得到性质 (iii)。

在本节其余部分，我们保持该命题的记号。

定理 2. --- 假设 \(u_{\mathrm{N}}\) 有有限迹。则函数 \(x \mapsto \mathrm{N}(x, x)\) 属于 \(\mathcal{L}^1(\mathrm{X})\) ，且我们有

\[
\mathrm{Tr} (u _ {\mathrm{N}}) = \int_ {\mathrm{X}} \mathrm{N} (x, x) d \mu (x).
\]

由条件 (i) 与 (ii)，我们有 \(|\mathrm{N}(x,x)| \leqslant \mathrm{H}(x,x)\) （对 \(x \in \widetilde{X}\)），故由条件 (iii)，函数 \(x \mapsto \mathrm{N}(x,x)\) 属于 \(\mathcal{L}^{1}(\mathrm{X})\) 。

条件 (i) 给出等式

\[
\mathrm{N} (x, x) = \sum_ {i \in \mathrm{I}} \alpha_ {i} \overline {{f _ {i} (x)}} g _ {i} (x)
\]

对所有 \(x \in \widetilde{\mathrm{X}}\) 成立。由条件 (ii) 与 (iii)，可以应用勒贝格定理（INT, IV, p. 137, § 3, no. 7，定理 6），由此可得

\[
\begin{array}{r} \int_ {\mathrm{X}} \mathrm{N} (x, x) d \mu (x) = \sum_ {i \in \mathrm{I}} \alpha_ {i} \int_ {\mathrm{X}} \overline {{f _ {i} (x)}} g _ {i} (x) d \mu (x) \\ = \sum_ {i \in \mathrm{I}} \alpha_ {i} \langle f _ {i} | g _ {i} \rangle = \mathrm{Tr} (u _ {\mathrm{N}}), \end{array}
\]

最后一个等式来自 IV, p. 170 命题 17 的推论。

引理 11. --- 若自同态 \(u_{\mathrm{N}}\) （\(\mathrm{L}^2(\mathrm{X})\) 的自同态）是正的，则 \(\mathrm{N}(x,x) \geqslant 0\) 对每个 \(x\) （\(\mu\) 的支撑中的元素）成立。

由于 \(u_{N}\) 是正的，可以选取族 \((f_{i})\) 与 \((g_{i})\) ，使得 \(f_{i} = g_{i}\) 对所有 \(i \in I\) 成立（参见 IV, p. 151 的注 3）。于是对每个 \(x \in \widetilde{X}\) ，我们有

\[
\mathrm{N} (x, x) = \sum_ {i \in \mathrm{I}} \alpha_ {i} | f _ {i} (x) | ^ {2} \geqslant 0.
\]

由于 \(N\) 连续且 \(\mu(\mathbf{X}-\widetilde{\mathbf{X}})=0\) ，函数 \(x \mapsto N(x,x)\) 在 \(\mu\) 的支撑上非负。

注. --- 该引理的逆命题不成立（IV, p. 316 的习题 14）。

命题 22. --- 假设 \(u_{\mathrm{N}}\) 是 \(\mathrm{L}^2(\mathrm{X})\) 的正自同态。则自同态 \(u_{\mathrm{N}}\) 有有限迹当且仅当函数 \(x \mapsto \mathrm{N}(x,x)\) 是 \(\mu\)-可积的。此时我们有

\[
\mathrm{Tr} (u _ {\mathrm{N}}) = \int_ {\mathrm{X}} \mathrm{N} (x, x) d \mu (x).
\]

若 \(u_{N}\) 有有限迹，定理 2 蕴含 \(x \mapsto \mathrm{N}(x, x)\) 在 X 上可积，且其积分是 \(u_{N}\) 的迹。

反之，设函数 \(x \mapsto \mathrm{N}(x, x)\) 可积。由于 \(u_{N}\) 是正的，可以选取标准正交族 \((f_i)_{i \in I}\) 与 \((g_i)_{i \in I}\) ，使得 \(f_i = g_i\) 对所有 \(i \in I\) 成立（IV, p. 151 的注 3）。对每个 \(x \in \widetilde{X}\) ，我们有

\[
\mathrm{N} (x, x) = \sum_ {i \in \mathrm{I}} \alpha_ {i} | f _ {i} (x) | ^ {2},
\]

由此，对 I 的每个有限子集 J，我们推出

\[
\begin{array}{c} \sum_ {i \in \mathrm{J}} \alpha_ {i} = \sum_ {i \in \mathrm{J}} \alpha_ {i} \int_ {\mathrm{X}} | f _ {i} (x) | ^ {2} d \mu (x) \\ = \int_ {\mathrm{X}} \sum_ {i \in \mathrm{J}} \alpha_ {i} | f _ {i} (x) | ^ {2} d \mu (x) \leqslant \int_ {\mathrm{X}} \mathrm{N} (x, x) d \mu (x). \end{array}
\]

因此族 \((\alpha_{i})_{i\in\mathrm{I}}\) 可和，这蕴含自同态 \(u_{N}\) 有有限迹（IV, p. 164 的命题 12）。

注. --- 即使当 X 紧且 N 连续时，自同态 \(u_{\mathrm{N}}\) （\(\mathrm{L}^{2}(\mathrm{X})\) 的自同态）也不总有有限迹（参见 IV, p. 314 的习题 8）。

